\documentclass[11pt,a4paper]{article}
\usepackage{amsmath,amssymb,amsthm,mathtools,booktabs,microtype}
\usepackage[margin=25mm]{geometry}
\usepackage[colorlinks=true,linkcolor=blue,citecolor=blue,urlcolor=blue]{hyperref}
\setlength{\emergencystretch}{2em}
\newtheorem{theorem}{Theorem}[section]
\newtheorem{lemma}[theorem]{Lemma}
\newtheorem{proposition}[theorem]{Proposition}
\newtheorem{corollary}[theorem]{Corollary}
\newcommand{\Q}{\mathbb Q}\newcommand{\Z}{\mathbb Z}
\DeclareMathOperator{\Tr}{Tr}\DeclareMathOperator{\Nm}{Nm}
\title{Amplified rationality descent:\\a cubic rational map and a sixteen-coordinate selector}
\author{\(\Z\Q\) Diophantine Project}
\date{19 November 2026}
\begin{document}\maketitle
\begin{abstract}
We audit the rationality-descent method of Pasten and Vidaux,
arXiv:2611.24748v1, and give several quantitative strengthenings.
Formal elliptic ratios approximate arbitrary rational targets
simultaneously at the real place and at finitely many finite places.
This supplies every rational input, eliminating the final quotient of
two separately defined sets. A two-sided odd-valuation transfer and an
odd denominator power amplify the resulting divisibility. The auxiliary
map can be $(T^3-2)/(T^3-3)$ instead of a polynomial of degree at least
seven. Prepared intrinsic pole equations suffice; their use gives a
definition of $\Q$ inside the ambient number field in sixteen total
coordinates. The assumptions are a specified truncated-height inequality,
positive elliptic rank stable from $\Q$, and the cited native-pole
completeness theorem. Ordinary $abc$ over the ambient field supplies the
height inequality. The correction for the finitely many irrational
exceptions is ineffective. These field-valued coordinate counts do not
reduce any of the separate counts for defining $\Z$ over $\Q$.
\end{abstract}

\section{Source audit and scope}

Pasten and Vidaux~\cite{PV} use an elliptic curve of positive rank
stable from $\Q$ to a number field $K$, two rational approximation
parameters, an even-pole predicate, and an $abc$ height inequality to
exclude all but finitely many irrational inputs. Their Theorem~1.1
assumes $abc$ over $K$ and
\[
 0<\operatorname{rank}E(\Q)=\operatorname{rank}E(K)
 \quad\text{for an elliptic curve }E/\Q.
\]
Their quadratic-field corollary uses Zywina's Theorem~1.2~\cite{Zywina}
with base field $\Q$. It does not supply this same stable $E/\Q$
for every higher-degree field. The source proof's use of a rational
approximation parameter at every prime over one rational prime, and its
passage to the Galois closure, are valid after excluding the fixed
ramified primes. We found no gap in that transfer under its stated
arithmetic inputs.

The descent architecture is due to \textbf{Hector Pasten and Xavier
Vidaux}. The results below are subsequent project strengthenings:
all-rational and real approximation, amplification, a two-sided cubic
map, a prepared-pole replacement, and a complete coordinate census.
The rational-integer polynomials used only in the final composition
corollary originate in \textbf{notes of Gabe Conant with Astra 6.2},
with the project audits and improvements recorded in v78--v81.

We do not independently prove $abc$, Zywina's theorem, or Burungale's
simultaneous-twist theorem. The latter is used in precisely its published
square classes through the native converse proved from it in
v80~\cite{Burungale,v80}. In particular, no square-class-one extension
of that theorem is assumed here.

\section{Approximation at all rational targets}

Fix an integral Weierstrass model of $E/\Q$ and an integer $m\ge1$
such that $\Gamma=mE(K)\subseteq E(\Q)$. Such an integer exists under
the rank hypothesis because $E(K)/E(\Q)$ is finite. Let
\[
 \xi=-x/y,\qquad
 \Gamma^\circ=\{P\in\Gamma\setminus\{O\}:x(P)y(P)\ne0\},
 \qquad
 \mathcal A=\{\xi(U)/\xi(V):U,V\in\Gamma^\circ\}.
\]
The set $\mathcal A$ is Diophantine over $K$ and is contained in
$\Q^\times$. Both membership statements follow from the algebraic
graph of $[m]$ and the explicitly guarded rational map $\xi$.

\begin{lemma}\label{lem:density}
Given a rational number $q$, finitely many rational primes, and prescribed
neighborhoods of $q$ in their completions and in $\mathbb R$, there
exists an element of $\mathcal A$ in all these neighborhoods.
It can also avoid any prescribed finite subset of $\Q$.
\end{lemma}
\begin{proof}
First let $q=a/b\ne0$, with nonzero integers $a,b$.
For a point $Q$ sufficiently deep in each formal group, put $z=\xi(Q)$.
The integral multiplication series give
\[
 \xi([a]Q)=az+z^2F_a(z),\qquad
 \xi([b]Q)=bz+z^2F_b(z),\qquad F_a,F_b\in\Z[[z]].
\]
If $v_p(z)>v_p(b)$, then
\[
 \frac{\xi([a]Q)}{\xi([b]Q)}-\frac ab
 =\frac{z\{bF_a(z)-aF_b(z)\}}{b\{b+zF_b(z)\}},
 \qquad
 v_p\left(\frac{\xi([a]Q)}{\xi([b]Q)}-\frac ab\right)
 \ge v_p(z)-2v_p(b).
\]
Thus an arbitrary prescribed finite-place accuracy follows by choosing
the formal neighborhoods sufficiently deep. Non-torsion points in these
neighborhoods and their nonzero multiples have $xy\ne0$; this follows
from the local expansions at $O$, not from non-torsion alone.

The subgroup of $\Gamma$ lying in all these neighborhoods has finite
index and contains a non-torsion point $R$. The compact group $E(\mathbb R)$
has nonzero multiples $Q_j=[n_j]R$ tending to $O$. They remain in all
the finite-place neighborhoods. The real formal expansion then gives
$\xi([a]Q_j)/\xi([b]Q_j)\longrightarrow a/b$. This proves the
simultaneous assertion for nonzero $q$.

For zero targets or finite avoidance, first choose a nonzero rational
$q'$ in the required open neighborhoods and outside the finite forbidden
set, by weak approximation. Shrink the neighborhoods around $q'$ to
exclude that set, and apply what was just proved. If the finite prime
set was empty, one may add a single prime to ensure the displayed
formal guards. No effective index bound is asserted.
\end{proof}

\section{Amplification and propagation of both valuation signs}

All valuations at finite primes of a number field are integral discrete
valuations; heights and counting functions are normalized by the field
degree. Let $f\in\Q(T)$ be nonconstant. For $g=f(t)\ne0,\infty$ and
$\lambda,\beta\in\mathcal A$ define
\begin{equation}\label{eq:inputs}
 A_g(\lambda)=\frac{(\lambda-g)^2(1+g)^2}{g^3},\qquad
 B_k(t,\lambda,\beta)
 =(\beta-t)^2\left(\frac{(1+\lambda)^2}{\lambda}\right)^{2k-1},
 \quad k\ge1.
\end{equation}
There is no division by $1+g$ or $1+\lambda$. Their zero fibers must
therefore not be silently removed.

Suppose a Diophantine pole predicate $\Psi$ excludes negative odd
valuations outside a fixed finite set of primes. Write $\Theta_{f,k}(t)$
for the existence of $\lambda,\beta\in\mathcal A$ such that both
expressions in~\eqref{eq:inputs} belong to $\Psi$, with the two displayed
nonzero denominators guarded.

\begin{lemma}\label{lem:local}
Let $p$ be a rational prime unramified in $K$, with every prime of $K$
above $p$ nonexceptional. Let $e=v_{\mathfrak p}(g)$ be nonzero and odd
at one prime $\mathfrak p$ above $p$.
If the first membership holds, then $v_{\mathfrak p}(\lambda)=e$.
At every prime $\mathfrak q$ above $p$,
both memberships imply
\[
             v_{\mathfrak q}(\beta-t)\ge k.
\]
\end{lemma}
\begin{proof}
If $\lambda=g$, the first conclusion is immediate. Otherwise its
valuation is
\[
 2v(\lambda-g)-3e+2\min(0,e),
\]
an odd integer. It must be positive. For $e>0$ this forces
$v(\lambda-g)>3e/2>e$; for $e<0$ it forces
$v(\lambda-g)>e/2>e$. In either case $v(\lambda)=e$.
As $\lambda\in\Q$, this valuation equals $e$ at every prime above
the same unramified rational prime.

At such a prime the valuation of $(1+\lambda)^2/\lambda$ is $-|e|$.
Unless $\beta=t$, the second input has odd valuation
\[
        2v(\beta-t)-(2k-1)|e|.
\]
It is positive, so $v(\beta-t)>(k-1/2)|e|\ge k-1/2$, proving
the integral lower bound. The case $\beta=t$ has infinite valuation
and already has rational input. If $\lambda=-1$, the first conclusion
excludes every consumed odd prime; no extra infinite exception is
created by the zero second input. The same observation treats $g=-1$.
\end{proof}

Let $D$ be an effective reduced divisor defined over $\Q$, supported
at zeros and poles of $f$ whose multiplicities are odd. Put $r=\deg D$.
At good reduction primes the odd contacts with $D$ are among the primes
of odd valuation of $f(t)$, including negative valuations. If $N_D^{\rm odd}$
counts those contacts once, then
\begin{equation}\label{eq:odd-lower}
 2N^{(1)}(D,t)\le N(D,t)+N_D^{\rm odd}(t)+O(1),\qquad
 N(D,t)\le r h(t)+O(1).
\end{equation}
Indeed, at every positive contact multiplicity $n$, the elementary
inequality $2\le n+(n\bmod2)$ applies. At a pole use a local reciprocal
parameter. Fixed bad primes contribute only bounded truncated errors.

\begin{theorem}\label{thm:amplification}
Assume, for all $t\in K$ outside $D$, the single-divisor estimate
\[
              N^{(1)}(D,t)\ge\theta h(t)-C,
              \qquad 2\theta-r>2/k.
\]
Then $\Theta_{f,k}(K)\setminus\Q$ is finite, independently of any
completeness assertion about $\Psi$.
\end{theorem}
\begin{proof}
Let $t$ be an irrational accepted input and choose an embedding $\sigma$
of $K$ into its Galois closure $L$ such that $\sigma(t)\ne t$.
Exclude the fixed rational primes ramified in $L$ or lying below a
pole-predicate exception. For each remaining rational prime $p$ below
an odd contact, Lemma~\ref{lem:local} gives
$v_{\mathfrak P}(\beta-t)\ge k$ at every $\mathfrak P\mid p$ in $L$.
Since $\beta$ is rational and $\sigma$ permutes these primes, it also
gives $v_{\mathfrak P}(\beta-\sigma(t))\ge k$. Hence
\[
 k\sum_{p\ {
m consumed}}\log p
 \le h(t-\sigma(t))\le2h(t)+\log2.
\]
The degree-normalized odd-contact count in $K$ is at most this sum
plus a fixed constant. Combining this bound with~\eqref{eq:odd-lower}
gives
\[
 (2\theta-r)h(t)-O(1)
 \le N_D^{\rm odd}(t)
 \le\frac{2}{k}h(t)+O(1).
\]
The strict gap bounds $h(t)$, and Northcott's theorem makes the
irrational accepted set finite. The pole or zero fibers excluded in
the definition are finite and do not affect this conclusion.
\end{proof}

Under $abc$ over $K$, its standard fixed-divisor consequence supplies
$\theta=r-2-\epsilon$. Thus $r-4-2\epsilon>2/k$ suffices.
The source's positive-contact polynomial construction can also be
amplified directly: replacing its second denominator $\lambda$ by
$\lambda^{2k-1}$ gives an upper bound $2h(t)/k+O(1)$.
A squarefree polynomial of degree five then suffices with $k=3$.
This increases a fixed exponent, without adding a witness coordinate.

\section{A cubic rational map}

Take
\begin{equation}\label{eq:cubic}
              f(T)=\frac{T^3-2}{T^3-3},\qquad k=2.
\end{equation}
The two cubics are squarefree and coprime: their discriminants are
$-108,-243$, and their difference is one. The divisor $D$ of their
six roots is reduced, rational, and has degree six. Its support comprises
six simple zeros or poles of $f$; infinity is neither. The three fibers
$f=0,\infty,-1$ contain no rational point, by the rational-root test
applied to $T^3-2,T^3-3,2T^3-5$.

Theorem~\ref{thm:amplification} needs only
\[
                  N^{(1)}(D,t)\ge(7/2+\eta)h(t)-C
                  \quad\text{for some }\eta>0.
\]
Ordinary $abc$ over $K$ gives this with any fixed
$0<\epsilon<1/2$, since $r-2-\epsilon=4-\epsilon$.
For example $\epsilon=1/4$ yields an odd-contact lower coefficient
$3/2$, an upper coefficient $1$, and a strict margin $1/2$.
Only this one fixed six-point divisor is needed. Increasing $k$ allows
any fixed lower coefficient $\theta>3$, with $k>1/(\theta-3)$.

\section{Prepared intrinsic poles suffice}

Let
\[
 \mathcal M=\{\pm1,\pm2,\pm3,\pm5,\pm6,\pm7,\pm10,\pm14\}.
\]
Define $\Psi(b)$ over any number field by the finite disjunction, for
$\mu\in\mathcal M$, of
\begin{equation}\label{eq:intrinsic}
 \exists v,s,Y:\quad
 \mu^5Y^2-\mu^2b(b+\mu v^2)^2s^4
                    +b(b-\mu v^2)^2(\mu-b)^2=0.
\end{equation}
Denote the polynomial on the left by $\Omega_\mu(b;v,s,Y)$.
Equivalently $Y^2=P_d(v,s)$, where $d=b/\mu$ and
\[
 P_d(v,s)=d(d+v^2)^2s^4-d(d-v^2)^2(1-d)^2.
\]
This is the intrinsic pole of v56~\cite{v56}; the clearing in
\eqref{eq:intrinsic} is literal, with no division by a variable.

\begin{lemma}\label{lem:prepared}
The predicate $\Psi$ excludes negative odd valuations outside primes
over $2,3,5,7$. It contains zero. It also contains every nonzero rational
$b$ with square reduced denominator such that
\[
 |b|<1,\quad v_2(b)\ge5,\quad
                v_3(b),v_5(b),v_7(b)\ge1.
\]
\end{lemma}
\begin{proof}
Outside the indicated primes $v(d)=v(b)$. Suppose this is negative
and odd. The valuations of $d$ and $v^2$ have opposite parity, so
$d+v^2$ and $d-v^2$ have the same finite valuation $a$.
The two terms of $P_d$ have respective valuations
\[
 v(d)+2a+4v(s),\qquad3v(d)+2a.
\]
If $s\ne0$ they are distinct odd integers, because equality would
give $2v(s)=v(d)$. Their difference cannot be a square. If $s=0$,
the second term alone has odd valuation. The case $v=0$ is covered
by $a=v(d)$. Thus no zero-coordinate loophole occurs. For $b=0$,
take $Y=0$.

For a prepared rational $b$, choose $\mu\in\mathcal M$ so that
$\mu b>0$ and $-\mu b\in\Q_2^{\times2}$. The sixteen multipliers
represent all real and two-adic square classes. Then $d=b/\mu$ is
positive and below one, is in $7\Q_2^{\times2}$, and has
$v_2(d)\ge4$. It still has square reduced denominator: the only
new denominators could be at $2,3,5,7$, but each prime occurs in $\mu$
with exponent at most one and the prescribed valuations prevent them.
The native converse of v80, derived from the published Theorem~B of
Burungale~\cite{Burungale,v80}, supplies the intrinsic witnesses.
Its equivalence with $P_d$ is established in v56. Alternatively, the
two curve points may be normalized and the first conic parametrized
as in the next section, with calibration equal to zero and $H=1-d$.
\end{proof}

\begin{proposition}\label{prop:rational-completeness}
With~\eqref{eq:cubic} and the predicate~\eqref{eq:intrinsic}, every
rational $t$ satisfies $\Theta_{f,2}(t)$.
\end{proposition}
\begin{proof}
Fix $t\in\Q$ and write $g=f(t)\in\Q\setminus\{0,-1\}$.
Choose $\lambda\in\mathcal A$ sufficiently close to $g$ at the real
place and at every prime in the support of $g(1+g)$, together with
$2,3,5,7$. Avoid $0,-1,g$. At every one of these finitely many primes,
the valuation of $A_g(\lambda)$ can be made arbitrarily large, and
its nonzero real absolute value arbitrarily small.
At every other prime, $g$ and $1+g$ are units, so
$v_p(A_g(\lambda))=2v_p(\lambda-g)$ is even. Consequently this first
input satisfies every preparation in Lemma~\ref{lem:prepared}.

Only now fix $\lambda$. Choose $\beta\in\mathcal A$ sufficiently
close to $t$ at the real place and at every prime in the supports of
$t$ and $\lambda(1+\lambda)$, together with $2,3,5,7$; omit the
support of $t$ if $t=0$. Avoid $t$. The nonzero second input in
\eqref{eq:inputs} is then arbitrarily small really and arbitrarily
integral at all these finite primes. Outside them its valuation is
$2v_p(\beta-t)$, hence even. It too satisfies the prepared-pole lemma.
This order includes every new numerator or denominator prime of the
actual chosen $\lambda$. No simultaneous prescribed square, or
unstated global approximation theorem, is being assumed.
\end{proof}

\section{Sixteen coordinates and finite corrections}

\begin{theorem}\label{thm:sixteen}
Suppose the stable-positive-rank hypothesis holds for $K$, the native
pole converse used in Lemma~\ref{lem:prepared} holds, and the six-point
height estimate of Section~4 holds. Then $\Q$ has a Diophantine
definition inside $K$ in at most \textbf{sixteen total $K$-coordinates},
including the input, with integer coefficients. In particular this
holds under $abc$ over $K$ and the cited arithmetic inputs.
\end{theorem}
\begin{proof}
The preceding results show that the set defined by $\Theta_{f,2}$
contains $\Q$ and has only finitely many irrational members.
Let $H\in\Z[T]$ be the product of the primitive minimal polynomials
of these exceptional members, or $H=1$ if there are none. It has no
rational root. Thus
\[
                         \Theta_{f,2}(t)\ \wedge\ H(t)\ne0
\]
defines $\Q$ exactly. The set, and therefore $H$, need not be effective.

For each of $\lambda$ and $\beta$, use two affine preimages on $E(K)$
under $[m]$. Four affine points use eight coordinates and four curve
equations. Their images under $[m]$, and then their two $\xi$-ratios,
are rational functions of these coordinates; no separate $\lambda$
or $\beta$ variable is needed. Standard division polynomials express
these maps on the required affine charts. Their nonzero denominators,
and the nonzero abscissa and ordinate of each image, are guarded.
Every desired point in $\Gamma^\circ$ has an affine preimage. Choose
the standard division-polynomial formulas for the coordinates of $[m]$:
their affine denominators vanish only at points mapping to the identity.
Consequently every preimage of a point in $\Gamma^\circ$ is permitted by
these charts. The additional nonzero image-coordinate guards hold by the
definition of $\Gamma^\circ$.

Each of the two intrinsic pole blocks uses precisely the three
coordinates $v,s,Y$. Substitute the rational expressions for their
inputs and clear denominators. One common coordinate $z_0$ imposes
$z_0\mathcal G=1$, where $\mathcal G$ is the product of every required
denominator, nonzero chart numerator, $f$ numerator, and $H(t)$.
There is no guard on $Y$, $s$, or $v$, and none on $1+\lambda$ is
needed. Thus the count is
\[
              1+8+2\cdot3+1=16.
\]
For each pole block, take the product of its sixteen cleared residuals
over $\mu\in\mathcal M$. The resulting equation asserts its finite
disjunction independently of the other block. Together with the four
curve equations and the common guard this gives seven polynomial
residuals. All coefficients are rational, and integral after clearing
fixed denominators. Each terminal residual~\eqref{eq:intrinsic} has
degree nine, thirteen monomials, and integer content one; across the
menu its maximum coefficient height is $14^5=537824$. The per-block
product multiplies the terminal degree by at most sixteen.

These seven equations may be combined without another coordinate.
Choose a rational prime $p$ unramified in $K$; the polynomial $Z^7-p$
is Eisenstein at any prime over $p$, so is irreducible over $K$.
The rational norm form of $\Q(p^{1/7})$ in its power basis is therefore
anisotropic after extension to $K$. Applying it to the seven residuals
is equivalent to their simultaneous vanishing. A product over a field
vanishes exactly when one factor vanishes, so the separate finite choices
introduce neither a coordinate nor a false branch. If $K$ has a real
embedding, the sum of the seven residual squares is an alternative
single equation; over a general number field a sum of squares alone
would not be valid.

For a finite degree bound, write the two substituted pole inputs as
$b_i=A_i/B_i$, with numerator and denominator degrees at most $d_0$.
Multiplication of $\Omega_\mu(A_i/B_i;v_i,s_i,Y_i)$ by $B_i^5$ gives
\[
\begin{split}
 &\mu^5Y_i^2B_i^5
 -\mu^2A_i(A_i+\mu v_i^2B_i)^2s_i^4B_i^2\\
 &\hspace{13mm}+A_i(A_i-\mu v_i^2B_i)^2(\mu B_i-A_i)^2.
\end{split}
\]
Its total degree is at most $5d_0+8$. If $g_0=\deg\mathcal G$, the
single norm equation therefore has degree at most
\[
                 7\max\{3,16(5d_0+8),g_0+1\}.
\]
This is a bound in the chosen multiplication maps and finite correction;
no uniform numerical degree or coefficient-height bound is asserted when
those data are ineffective.
\end{proof}

There is an alternative finite correction with a degree bound independent
of the cardinality of the exceptional set. The power-shift lemma of
v69~\cite{v69} supplies a fixed integer $B$ avoiding that finite set
on the image of $x\mapsto x^\ell+B$, where $\ell$ reflects rationality
in $K$. Then $\Theta_{f,2}(x^\ell+B)$ defines $\Q$, in the same sixteen
coordinates. Its degree is at most $\ell$ times that of the uncorrected
polynomial. One may take any prime $\ell$ with $\ell-1>[K:\Q]$.
The shift can lie in a prescribed integral progression, but is generally
ineffective. This construction also preserves rational coefficients.

\begin{corollary}\label{cor:composition}
If a polynomial in $N$ total rational coordinates defines $\Z$ over
$\Q$, then the selector above gives a definition of $\Z$ inside $K$
in $16N$ total $K$-coordinates. The ring of integers of $K$, of degree
$d$ over $\Q$, has a definition in $16dN$ total $K$-coordinates.
Using the project's eleven-coordinate polynomial gives $176$ and
$176d$, respectively, under that polynomial's explicitly retained
arithmetic inputs as well as the hypotheses above.
\end{corollary}
\begin{proof}
Restrict each of the $N$ original coordinates, including its input,
using an independent copy of the sixteen-coordinate selector. The
total is $N+15N=16N$. For the ring, use an integral basis and eliminate
the last basis coefficient by the distinguished input, as in v12
and v69~\cite{v12,v69}. This gives exactly $d$ times the previous
coordinate count. Restricting the original input as well as the
witnesses is essential. Finally $abc$ over $K$ implies ordinary
rational $abc$ by restricting to rational triples; outside finitely
many ramified primes the normalized radical contributions agree.
\end{proof}

For a quadratic field, Zywina's cited theorem supplies the stable
rank-one curve, so the corresponding counts are $176$ for $\Z$ and
$352$ for its ring of integers, with the hypotheses in the corollary.
These are existence bounds over that quadratic field. They are not
rational-variable counts for a polynomial over $\Q$, nor quartic-degree
claims, nor effective specifications of a sufficient correction parameter.

\section{An explicit full Pheidas menu}

Although the sixteen-coordinate construction needs only prepared inputs,
the full even-pole predicate in the source can also be made explicit.
For $\mu\in\mathcal M$ choose $\gamma$ from the following sets and put
$\delta=\mu\gamma$:
\[
\begin{array}{c|c}
\mu&\gamma\\\hline
\pm1,\pm2&7\\
\pm3,\pm6&7,23,39,15\\
\pm5,\pm10&31,7,55,15\\
\pm7,\pm14&15,31,7,119
\end{array}
\]
There are $4+16+16+16=52$ choices. Every $\gamma$ is positive and
congruent to seven modulo eight. For each $p=3,5,7$, its four listed
values represent the four square classes of $\Q_p^\times$, as checked
by valuation parity and the Legendre symbol of the unit part.
Consequently, after $\mu b>0$ and $-\mu b\in\Q_2^{\times2}$ have
been arranged, one can choose $\gamma$ so that $\delta b$ is positive
and a square at every prime dividing $2\mu$. These are exactly the
local conditions needed in Burungale's completeness proof; its remaining
finite support conditions are imposed on the free calibration parameter.
All excluded primes lie in
\[
                     \{2,3,5,7,11,13,17,23,31\}.
\]

For completeness, the two curve conditions of that full predicate admit
the following smaller chart. Write $a=\mu b$, let $t_0$ be its calibration
parameter, put
\[
 U=\frac{1-\delta b t_0^2}{\mu(1+\delta b t_0^2)},\qquad H_0=1-aU^2,
\]
and normalize the first square-abscissa multiplier to one. The resulting
conics are $Y_0^2=a^3-aE^2$ and
$Z_0^2=a^3w^4-aE^2H_0^2$. Their first conic is parametrized by
\[
 E=\frac{a(a-r^2)}{a+r^2},\qquad
 Y_0=\frac{2a^2r}{a+r^2},
\]
and the second becomes
\[
 Z_1^2=a(a+r^2)^2w^4-a(a-r^2)^2H_0^2.
\]
This uses only $t_0,r,w,Z_1$, with the displayed denominators and
$a(a-r^2)wH_0$ nonzero. The missing first-conic point $E=-a$ is
irrelevant for completeness: the supplied positive twist parameter
and positive first abscissa give $E>0,a>0$. The zero input is a separate
trivial branch. This chart supplies an eighteen-coordinate alternative
to Theorem~\ref{thm:sixteen} without the real-preparation simplification;
it is not the preferred sixteen-coordinate result.

\section{Limits and verification}

The decisive field comparison is $t-\sigma(t)\ne0$. It disappears
when $t\in\Q$, so this argument does not distinguish rational integers
from other rationals. Neither formal approximation nor real recurrence
imposes a new exact square at a prescribed original detector tuple.
In particular, the separate $25/11$ constructions still need their
global rational-realization arguments for a variable saving.

The independent checker verifies the cubic divisors, exact margins,
the cleared intrinsic polynomial, local valuation cases, conic maps,
the multiplier menus, and coordinate census. Universal approximation,
valuation propagation, finite correction, and composition are proved
above and reviewed separately. A finite valuation grid is only a
regression check. No computation here establishes the cited arithmetic
theorems or an effective exceptional-height bound.

\begin{thebibliography}{9}\footnotesize\raggedright
\setlength{\itemsep}{1pt}
\bibitem{PV} H.~Pasten and X.~Vidaux,
\emph{On Hilbert's tenth problem for number fields},
arXiv:2611.24748v1, 18 November 2026,
\url{https://arxiv.org/html/2611.24748v1}.
\bibitem{Zywina} D.~Zywina,
\emph{Rank one elliptic curves and rank stability},
arXiv:2505.16960v1, Theorem~1.2,
\url{https://arxiv.org/html/2505.16960v1}.
\bibitem{Burungale} A.~Burungale,
\emph{A Diophantine definition of $\Z$ in $\Q$},
arXiv:2611.16130v1, Theorem~B and Section~3,
\url{https://arxiv.org/html/2611.16130v1}.
\bibitem{v80} \(\Z\Q\) Diophantine Project,
\emph{One-helper pole normalization: an ordinary-abc quartic in 58 coordinates},
\nolinkurl{papers/quartic_existence/one_helper58_20261119_v80.tex}.
\bibitem{v56} \(\Z\Q\) Diophantine Project,
\emph{Intrinsic pole predicate},
\nolinkurl{papers/pole_parameters/intrinsic_pole_20261117_v56.tex}.
\bibitem{v69} \(\Z\Q\) Diophantine Project,
\emph{Power shifts and finite exceptions in Diophantine definitions},
\nolinkurl{papers/number_fields/power_shifts_and_finite_exceptions_20261117_v69.tex}.
\bibitem{v12} \(\Z\Q\) Diophantine Project,
\emph{Number-field definability and quantitative composition},
\nolinkurl{papers/number_fields/number_field_definability_20261109_v12.tex}.
\bibitem{v78} G.~Conant, with Astra 6.2, supplied construction notes;
project audit and improvements,
\nolinkurl{papers/conditional_few_variables/twelve_contact_first_scale_20261119_v78.tex}.
\end{thebibliography}
\end{document}
